% !PRINTED none
% No running head/page number visible anywhere on this scan: an Abschnitt-opening
% page (title + rule) with the § 57 heading and its opening paragraphs and formulas
% (1),(2) beginning directly beneath it, all on the same physical page.

\begin{center}
\gesperrt{Sechster Abschnitt.}

\textbf{Multiplikation und Teilung der elliptischen Funktionen.}
\end{center}

\bigskip
\begin{center}
\rule{2cm}{0.4pt}
\end{center}
\bigskip

\begin{center}
\textbf{§ 57. \gesperrt{Multiplikation der elliptischen Funktionen.}}
\end{center}

Unter der Multiplikation der elliptischen Funktionen versteht man
die Darstellung der Funktionen $\mathrm{sn}\,nv$, $\mathrm{cn}\,nv$,
$\mathrm{dn}\,nv$ für ein ganzzahliges $n$ als rationale Funktionen
von $\mathrm{sn}\,v$, $\mathrm{cn}\,v$, $\mathrm{dn}\,v$, eine
Aufgabe, die, wie aus dem Additionstheorem ersichtlich ist, immer
gelöst werden kann.

Die Form der Lösung ergibt sich leicht aus der Betrachtung der
$\vartheta$-Funktionen.

Es sind, wie aus § 21 unmittelbar zu ersehen, die Funktionen
$\vartheta_{g_1,g_2}(nu)$ $\Theta$-Funktionen der Ordnung $n^2$,
deren Charakteristik bei geradem $n$ gleich $(0,0)$, bei ungeradem
$n$ gleich $(g_1,g_2)$ ist. Überdies ist $\vartheta_{11}(nu)$ eine
ungerade, $\vartheta_{00}(nu)$, $\vartheta_{10}(nu)$,
$\vartheta_{01}(nu)$ sind gerade Funktionen von $u$. Es lassen sich
also diese Funktionen rational durch die $\vartheta$-Funktionen
darstellen und die Sätze des § 21 ergeben die Form dieser
Ausdrücke. Es wird, wenn wir wieder mit $F^{(\nu)}(x,y)$ eine
ganze, rationale, homogene Funktion $\nu$ter Ordnung bezeichnen:

\gesperrt{bei geradem $n$:}

\weq{1}{\begin{aligned}
\vartheta_{11}(nu) &= \vartheta_{00}(u)\vartheta_{01}(u)\vartheta_{10}(u)\vartheta_{11}(u)\,F^{\frac{n^2-4}{2}}\bigl[\vartheta_{11}^2(u),\vartheta_{01}^2(u)\bigr],\\
\vartheta_{g_1,g_2}(nu) &= F^{\frac{n^2}{2}}\bigl[\vartheta_{11}^2(u),\vartheta_{01}^2(u)\bigr], \quad (g_1,g_2)=(10),(01),(00),
\end{aligned}}

\gesperrt{bei ungeradem $n$:}

\weq{2}{\vartheta_{g_1,g_2}(nu) = \vartheta_{g_1,g_2}(u)\,F^{\frac{n^2-1}{2}}\bigl[\vartheta_{11}^2(u),\vartheta_{01}^2(u)\bigr].}
